% Lecture visuelle directe avec la vision du modèle ; aucun OCR.
% Source : PDF original, vue 064, côté gauche, page manuscrite 128.
% Statut : transcrit ; relecture indépendante à effectuer.
% Chapitre complet : p. 128–145 (début).
\chapter{Préfecture de Monsieur Chaulet}
\begin{center}
Établissement des confréries du St Sacrement. Bref d'indulgences. Établissement de la fête d'actions de grâces. Association aux diverses Congrégations.
\end{center}

\section*{Son installation}
% Mention marginale de la page manuscrite 128.
On installa le nouveau préfet conformément aux prescriptions du règlement le jour de la fête de la Nativité de la Ste Vierge. La Congrégation a toujours eu pour principe d'entourer celui qu'elle a revêtu de la présidence d'un certain décorum afin que cela contribuât à le faire respecter et mieux obéir de tous. C'est dans cette vue qu'elle établit pour l'installation de M. Chaulet les mêmes formes et cérémonies qui depuis lors sont restées constamment en usage parmi nous.

Notre nouveau préfet était un de ces anciens élèves du séminaire de St Irénée dont nous avons parlé dans le Précis historique qui a précédé cette histoire. Il avait été reçu l'année précédente à la même fête de la Nativité de la Ste Vierge ; il avait 32 ans lorsqu'il entra dans la Congrégation. M. Chaulet avait fait d'excellentes études ; quoique par état il fût obligé de s'occuper de commerce il avait du goût pour la littérature et surtout pour la poésie ; il se plaisait beaucoup à faire usage de ses talents dans l'intérêt de la religion ; nous avons déjà vu qu'il saisissait avec empressement l'occasion de les rendre utiles à la Congrégation. Il était naturellement grave ; son allure était un peu empesée ;
% Source : PDF original, vue 064, côté droit, page manuscrite 129.
elle aurait pu paraître affectée à ceux qui l'auraient moins connu ; mais au fond personne ne pouvait être meilleur, plus simple, plus dépouillé de toute espèce de prétentions. La Congrégation était alors composée en partie de bien jeunes gens ; il savait bien s'en faire respecter. Les discours qu'il préparait avec soin pour les adresser aux futurs congréganistes lorsqu'ils étaient présentés en imposaient à ces nouveaux reçus et leur imprimaient un grand respect pour la Congrégation. C'était donc dans l'état des choses précisément le préfet qui nous convenait.

Il n'occupait pas depuis bien longtemps le fauteuil de la présidence lorsque la Providence vint employer pour la première fois d'une manière directe la Congrégation à procurer sa gloire. Les congréganistes avaient observé avec peine à l'époque où l'exercice du culte extérieur s'était rétabli dans notre ville par les processions du St Sacrement qu'un grand nombre de chrétiens cependant bien sincèrement attachés à la religion s'étaient dispensés d'assister aux processions et se contentaient de les regarder passer. Ils laissaient à de pauvres artisans ou à des domestiques la fonction honorable de prendre une part directe au triomphe de J.-C. Il n'y avait certainement pas de l'irréligion dans cette indifférence apparente ; c'était un peu d'apathie, peut-être une vieille habitude car il n'était que trop vrai que cela se pratiquait ainsi avant la Révolution. Il serait bien possible qu'il s'y mêlât aussi un peu de respect humain. Quoi qu'il en soit Dieu inspira à la Congrégation la bonne pensée de chercher à y porter remède. Ses chefs crurent que l'établissement dans les paroisses de notre ville des confréries
% Source : PDF original, vue 065, côté gauche, page manuscrite 130.
du St Sacrement en observant toutefois d'y faire entrer surtout les personnes notables de chaque paroisse dans l'heureuse habitude de remplir leurs devoirs de religion suffirait pour amener toutes ces personnes à se présenter à l'avenir pour former le cortège de Notre Seigneur dans les cérémonies extérieures. La Congrégation pensait trouver encore dans cette création de confréries un avantage personnel ; elle pourrait regarder les jeunes gens qui y seraient admis comme formant pour nous une espèce de noviciat où nous pourrions choisir ceux qui nous paraîtraient propres à grossir nos rangs ; en un mot c'était former en même temps une pépinière pour la Congrégation.

\section*{Confréries du St Sacrement}
% Mention marginale de la page manuscrite 130.
L'expédient une fois trouvé et regardé comme avantageux sous tous les rapports il ne restait plus qu'à chercher les moyens d'exécution. La première chose à obtenir c'était l'agrément de MM. les curés car il était évident que dans une affaire de ce genre sans leur intervention il devenait absolument impossible de rien entreprendre. La Congrégation s'étendant dans toute la ville avait des membres dans presque toutes les paroisses mais cependant les congréganistes n'étaient encore un peu nombreux que dans celles de St Nizier, de St Pierre et de St François de Sales. Ce fut donc dans ces trois paroisses qu'elle se détermina à essayer l'exécution de son plan. Elle délégua dans chacune d'elles un des paroissiens choisi parmi ses membres qu'elle chargea de soumettre à son pasteur en son propre et privé nom le projet du nouvel établissement. Cette tentative eut un succès complet. MM. les curés n'eurent
% Source : PDF original, vue 065, côté droit, page manuscrite 131.
pas de peine à comprendre quels avantages ils pourraient trouver dans ces confréries soit pour augmenter la splendeur du culte divin soit pour procurer le bien spirituel de leurs paroissiens ; ils s'empressèrent donc de favoriser leur établissement. Lorsque le moment arriva de célébrer de nouveau la Fête-Dieu en 1804 on vit avec consolation que non seulement la plupart des notables paroissiens de ces trois paroisses s'étaient très volontiers décidés à faire partie des nouvelles confréries mais que leur exactitude à remplir les devoirs de confrères et surtout à assister aux processions faisait l'édification de tous les paroissiens. La Congrégation a travaillé plus tard à former ces confréries dans les paroisses où ses membres se sont multipliés comme par exemple à St Polycarpe. Mais comme le bien est contagieux aussi que le mal une fois l'impulsion donnée toutes les autres paroisses ont peu à peu suivi l'exemple qui leur avait été offert en sorte que depuis bien longtemps il n'y a pas une seule paroisse de Lyon qui n'ait sa confrérie du St Sacrement.

La Congrégation y a trouvé un moyen constant de procurer la gloire de Dieu. Ces sociétés formées par elle aux règlements desquelles les congréganistes ont travaillé et dont souvent ils ont été les principaux rédacteurs et qui dans leurs bureaux comptent un grand nombre de congréganistes ont puisé à la source d'où elles sont sorties l'esprit qui les dirige et cet esprit s'est conservé sauf toutefois les modifications qu'exige le but spécial pour lequel elles ont été établies. On peut dire que les confréries du St Sacrement sont entrées dans la sphère que nous verrons s'aggrandir et se développer plus tard des œuvres dirigées par la Congrégation.
% Source : PDF original, vue 066, côté gauche, page manuscrite 132.
Cet établissement a produit encore un autre avantage : il a mis un grand nombre de congréganistes en rapport immédiat avec leurs curés. Ceux-ci les ont trouvés toujours disposés à seconder de tout leur pouvoir tous les projets que le zèle des pasteurs pouvait former pour l'avantage de leurs paroisses. Ceux d'entre ces vénérables curés qui sur une connaissance superficielle de la Congrégation auraient pu craindre que la direction qui lui était imprimée par un corps religieux n'éloignât de leurs paroisses les personnes qui en font partie purent se convaincre qu'ils trouvaient au contraire dans les congréganistes les paroissiens les plus dévoués et les plus exacts à remplir les devoirs particuliers que l'Église nous impose en cette qualité.

\section*{Bref d'indulgences}
% Mention marginale de la page manuscrite 132.
Cette nouvelle création augmentait déjà sensiblement l'importance que la Congrégation avait acquise. Une circonstance nouvelle vint à peu près en même temps lui donner encore plus de consistance. Nous étions alors trop peu de chose, nous ne comptions pas assez dans le nombre des établissements utiles à l'Église pour que nous eussions osé nous adresser à Rome pour en obtenir des indulgences. Mais ce que nous n'aurions pas cru pouvoir faire notre père le tenta en notre faveur. Le cardinal légat était encore à Paris ; il avait reçu du pape les pouvoirs les plus étendus pour accorder au moins temporairement ces faveurs spirituelles en son nom. Le P. Roger rédigea une supplique et la fit recommander puissamment. La demande eut pour résultat un succès complet. Nous reçûmes un rescrit du cardinal légat par lequel il nous accordait pour sept
% Source : PDF original, vue 066, côté droit, page manuscrite 133.
ans des indulgences sans doute moins abondantes que celles dont nous jouissons maintenant mais pourtant bien au-dessus de ce que nous espérions. À cette heureuse nouvelle la Congrégation tressaillit d'allégresse ; tout nous réussissait alors au gré de nos désirs. Chaque membre repassait avec émotion dans son cœur toutes ces circonstances que la Providence amenait tour à tour pour nous confirmer dans la jouissance de ces biens spirituels dont elle avait daigné nous mettre en possession. Comment ne pas être vivement attendris en considérant tous ces moyens de sanctification amassés avec profusion autour de nous pour nous rendre le salut facile. Ces indulgences qui venaient de nous être adressées au nom du souverain pontife nous semblaient apporter avec elles la sanction donnée par l'Église à l'établissement de notre chère Congrégation ; c'était la clef de voûte qui venait consolider tout l'édifice.

\section*{Fête d'actions de grâces}
% Mention marginale de la page manuscrite 133.
Une pensée de gratitude devait être pour tous la suite naturelle de ces réflexions ; aussi ce sentiment remplissait-il nos âmes ; ce n'est pas assez dire, il surabondait ; nos cœurs ne pouvaient plus le contenir ; il fallait enfin qu'il se fît jour, qu'il éclatât au-dehors. Pour répondre à cet élan général il était indispensable de donner à notre reconnaissance une manifestation extérieure et durable. C'est cette disposition de tous les cœurs congréganistes qui détermina la Congrégation à établir une fête solennelle d'actions de grâces. Elle perpétuera par là le souvenir des bienfaits que la main miséricordieuse de notre Dieu avait répandus si abondamment sur nous. Elle perpétuait aussi le souvenir de l'impression profonde
% Source : PDF original, vue 067, côté gauche, page manuscrite 134.
qu'ils avaient faite sur nos cœurs et pour qu'on n'oubliât jamais que la réception des indulgences avait le premier rang parmi les grâces dont elle se proposait de remercier le Seigneur la Congrégation voulut que le bref qui nous en mettait en possession fût lu chaque année en ce jour de fête. Cette solennité fut dès lors fixée au quatrième dimanche après Pâques ainsi qu'elle se célèbre encore aujourd'hui.

\section*{Association avec la Congrégation de Paris}
% Mention marginale de la page manuscrite 134.
Les rapports de la Congrégation venaient de s'aggrandir dans l'intérieur de la ville par son influence sur les confréries du St Sacrement qu'elle avait fondées. Nous allons les voir maintenant s'aggrandir à l'extérieur par des associations avec d'autres Congrégations.

À peu près dans le même temps que le Seigneur avait inspiré à quelques jeunes Lyonnais la pensée de se réunir en Congrégation il avait suscité à Paris un saint prêtre, reste vénérable de l'ordre des jésuites, et il l'avait chargé de former aussi dans cette grande ville une Congrégation en l'honneur de la Ste Vierge. Le Père Delpuy (c'était le nom de cet homme de Dieu) travailla donc aussitôt à réunir autour de lui tous les restes d'une jeunesse demeurée fidèle au culte du vrai Dieu et dans une ville où de grandes vertus brillent d'un éclat d'autant plus vif qu'elles se montrent au centre de l'impiété et de la corruption. Cet appel fut entendu ; la nouvelle Congrégation compta bientôt au nombre de ses membres des hommes les plus distingués soit par leur naissance soit par leurs talents et cette société fut formée sur le modèle des anciennes Congrégations des jésuites. Elle était toute entière dans
% Source : PDF original, vue 067, côté droit, page manuscrite 135.
les mains de son directeur qui décidait à peu près seul de toutes les réceptions. Le but de la Congrégation, but qui était secondé par tout le zèle du directeur, était la sanctification de tous ceux qui la composaient ; sa direction tendait à la communion fréquente et pour les y animer davantage une réunion générale avait lieu tous les quinze jours dans une chapelle particulière. Après une lecture faite par un congréganiste le directeur célébrait la messe. La plupart des congréganistes présents s'approchait de la table sainte ; on procédait aux réceptions. Après la messe on donnait lecture des lettres des membres absents. Il n'y avait ni réunions particulières pour les bonnes œuvres ni aucun de ces exercices si variés qui font la vie de la Congrégation de Lyon et réchauffent sans cesse la charité dans le cœur des congréganistes. Vous voyez par là que la Congrégation de Paris avait un but beaucoup plus restreint que la nôtre et que l'esprit de prosélytisme n'y entrait pour rien. Mais à cette époque cette société était dans toute sa ferveur. Elle comptait parmi ses membres des saints et en grand nombre. Ce but que leur société n'avait pas d'une manière directe Dieu l'avait placé dans leurs cœurs. Ils ne perdaient donc aucune occasion d'étendre le règne de la foi. La nouvelle de l'existence d'une Congrégation à Lyon leur parvint d'abord d'une manière assez vague mais ils ne tardèrent pas à l'éclaircir. L'union des Congrégations leur parut pouvoir seconder merveilleusement la cause de la foi ; ils cherchèrent donc des moyens de se mettre en rapport avec nous. M. le duc de Montmorency eut sur ces entrefaites un motif pour faire un voyage à Lyon. Il fut délégué par la Congrégation pour chercher à nous connaître et s'il y avait lieu conclure une
% Source : PDF original, vue 068, côté gauche, page manuscrite 136.
alliance entre les deux Congrégations. Pour se faciliter les moyens de se mettre en rapport avec la Congrégation de Lyon M. de Montmorency se procura des lettres de recommandation par l'intermédiaire de M. Jordan l'un de nos fondateurs qui faisait alors son séminaire à Paris.

Aussitôt que ce digne et respectable représentant de la Congrégation de Paris fut arrivé à Lyon il s'empressa de nous chercher. Il nous dit avec les marques de la plus touchante bonté qu'il considérait cette grande affaire comme le but principal de son voyage. La Congrégation de Lyon n'eut pas de peine à comprendre quel avantage elle pouvait retirer de cette union soit dans l'intérêt du bien qu'elle se proposait d'opérer pour lequel elle trouverait dans bien des occasions appui et protection soit pour l'utilité particulière de ceux de nos congréganistes que leurs affaires ou leurs études appelleraient à Paris. Le préfet désigna donc un plénipotentiaire pour rédiger un traité d'alliance de concert avec le plénipotentiaire de la Congrégation de Paris. Tout fut promptement d'accord : il fut convenu qu'il y aurait union de prières entre les deux Congrégations et qu'elles recevraient réciproquement dans leurs assemblées ceux de leurs membres qui passeraient ou séjourneraient dans leurs villes respectives sur une simple lettre de recommandation du préfet. Ainsi outre les moyens que nos congréganistes devaient trouver à Paris pour résister à la contagion de la capitale cette association allait offrir à tous sans exception un sujet sans cesse renouvellé d'édification par le bon exemple qu'ils recevraient des congréganistes de Paris que leurs affaires
% Source : PDF original, vue 068, côté droit, page manuscrite 137.
ou leurs plaisirs amèneraient à Lyon et ceux-ci de leur côté devaient trouver bien doux de se voir accueillis comme des frères par une société qui devenait leur famille. Aussi cette association contribua alors de toutes manières à fortifier notre Congrégation ; elle en apprécia si bien les avantages qu'elle voulut rendre plus intimes encore les relations qui nous unissaient déjà avec quelques-uns de ces hommes vertueux si distingués par leurs éminentes qualités. Elle prit un arrêté par lequel elle se donnait la faculté d'associer personnellement à la Congrégation de Lyon ceux des membres des Congrégations associées auxquels elle jugerait à propos d'accorder ce titre. Ces associés ont joui dans la Congrégation à peu près des mêmes droits dont jouissent actuellement les membres honoraires avec lesquels comme nous le verrons plus tard ils sont maintenant confondus. Nous avons associé ainsi un assez grand nombre de membres de la Congrégation de Paris ; je n'en citerai que quelques-uns.

\section*{M. Matthieu de Montmorency}
% Mention marginale de la page manuscrite 137.
Parmi eux nous placerons d'abord et au premier rang l'illustre plénipotentiaire qui avait bien voulu se charger lui-même de nous offrir une alliance qui nous était si avantageuse. Ce fut un bien grand bonheur, Messieurs, pour notre Congrégation alors encore si foible de pouvoir considérer comme lui appartenant M. le vicomte de Montmorency si digne du nom glorieux de premier baron chrétien, de pouvoir regarder comme un frère cet homme bien plus recommandable encore par sa grande vertu que par sa grande naissance. Nous ne craignons pas d'être démentis en avançant ici que M. de Montmorency s'est montré constamment un parfait modèle de charité, de douceur, d'humilité tout évangélique. Il était très jeune au moment
% Source : PDF original, vue 069, côté gauche, page manuscrite 138.
où commencèrent nos troubles politiques. Il avait le malheur d'avoir eu pour précepteur cet abbé Sieyès qui fut un des coryphées de notre Révolution. Les conseils perfides d'un guide aussi dangereux égarèrent sa belle âme ; il fit alors quelques faux pas par excès de candeur et de vertu ; mais comme il s'est glorieusement relevé d'une chute involontaire que son inexpérience seule avait causée ! Comme sa vie entière a été employée à la réparer par la plus noble conduite ! Déjà à l'époque où il vint à Lyon pour signer avec nous le premier traité diplomatique auquel il ait apposé sa signature il avait mérité par le zèle qu'il déployait pour favoriser l'établissement de toutes les œuvres formées dans l'intérêt de l'humanité souffrante d'être universellement appelé le ministre des bonnes œuvres et ce beau ministère dans quelque position qu'il se soit trouvé il l'a rempli toute sa vie. M. de Montmorency était porteur d'un nom qui ne le cède en noblesse qu'à celui de la famille de nos rois ; il était rempli de connaissances et de talents. Avec tant d'avantages qui le distinguaient du commun des hommes il eût été difficile de se montrer plus modeste. On se sentait pénétré d'admiration et attendri jusqu'aux larmes lorsqu'on le voyait et cela en toutes circonstances saisir avec cette simplicité et cette candeur qui ne se rencontrent que dans les hommes de foi toutes les occasions de déplorer amèrement les fautes de sa jeunesse. Que de fois lorsqu'on jetait devant lui quelque blâme sur la conduite de quelqu'un ne lui ai-je pas vu prendre aussitôt en main la cause de la personne accusée et tout en cherchant
% Source : PDF original, vue 069, côté droit, page manuscrite 139.
à la défendre rentrer en même temps en lui-même par une pensée d'humilité et laisser échapper de son cœur ces mots tout empreints d'un sentiment de contrition : Eh ! qui dans la Révolution n'a pas eu quelque chose à se reprocher et moi surtout ?

Laissez-moi vous ajouter encore que de toutes les opinions qu'il a prononcées à la tribune soit en qualité de pair soit en qualité de ministre il n'en a envoyé qu'une à l'auteur de cette histoire qu'il honorait de son amitié et c'est celle où en présence de la France et à la face du monde entier il déplorait et condamnait ouvertement les égarements de ses premières années.

M. de Montmorency après avoir parcouru en saint soit sa carrière privée soit sa carrière publique est mort aux pieds de la croix un Vendredi saint à l'heure même de la mort du Sauveur. Digne fin qui a couronné la plus belle vie par la mort la plus sainte ! Jésus-Christ se rendant obéissant jusqu'à la mort et à la mort de la croix avait rempli son cœur de cette humilité et de cette douceur qu'il est venu faire connaître à la terre. Il s'est chargé de récompenser la fidélité de son serviteur à suivre ses leçons et ses exemples : \textit{Discite a me quia mitis sum et humilis corde}\footnote{Traduction : « Apprenez de moi, car je suis doux et humble de cœur. » Référence : Matthieu 11, 29. (Note du transcripteur).}.

\section*{M. Alexis de Noailles}
% Mention marginale de la page manuscrite 140 ; placée au début du portrait.
Il y avait à peine quelques semaines que l'association avec la Congrégation de Paris avait été conclue lorsque nous reçûmes la visite du comte Alexis de Noailles. M. de Noailles était alors dans toute l'ardeur de la première jeunesse. Son zèle bouillant pour la cause de la foi vint animer et réchauffer le nôtre. Il assista à plusieurs de nos assemblées et nous lut un rapport sur diverses sociétés de bonnes œuvres,
% Source : PDF original, vue 070, côté gauche, page manuscrite 140.
qu'il avait visitées dans le cours d'un voyage qu'il venait de faire dans le midi. Le caractère de M. de Noailles se trouvait singulièrement en rapport avec l'esprit qui a toujours régné dans la Congrégation de Lyon. Avec l'activité naturelle dont il était doué et qui se portait toute entière sur ce qui pouvait contribuer au triomphe de la religion il ne pouvait manquer d'apprécier cette action sans cesse renouvellée avec laquelle notre Congrégation entretient et ranime constamment la vie de la foi dans le cœur de tous ses membres. Il se présenta à nous avec des manières si franches et si ouvertes qu'il eut bientôt gagné tous nos cœurs. Aussi notre liaison avec lui devint-elle très intime. Il se considérait tout à fait comme un congréganiste de Lyon et il avait conçu la plus haute estime pour la Congrégation. Il se plaisait à dire par manière de plaisanterie qu'il classait ainsi les trois préfets dont Lyon était la résidence : celui qui sans contredit occupait le premier rang était celui de la Congrégation des jeunes gens ; venait ensuite le préfet de la Congrégation des hommes et enfin il voulait bien par grâce admettre en troisième lieu le préfet du département du Rhône. On peut dire aussi de ce badinage qu'il établissait dans toute sa réalité l'ordre de ses affections. Nous retrouverons encore plusieurs fois M. de Noailles dans la suite de cette histoire.

\section*{M. de Rohan}
% Mention marginale de la page manuscrite 140.
M. de Rohan-Chabot était encore dans la première jeunesse (je crois qu'il n'avait guère que quinze ans) lorsque la Congrégation le reçut en qualité d'associé ; il lui fut présenté par M. de Noailles. Dans une vie
% Source : PDF original, vue 070, côté droit, page manuscrite 141.
variée de bien des manières différentes il a constamment persévéré dans la pratique de la vertu. Il a fini par entrer dans le saint ministère et est devenu archevêque de Besançon. Il est mort il y a peu d'années dans sa ville archiépiscopale après avoir édifié ses diocésains soit par son zèle apostolique soit par la tendre dévotion qu'il a toujours eue envers la Ste Vierge.

\section*{M. Teyssier}
% Mention marginale de la page manuscrite 141.
Nous avons encore dû à la Congrégation de Paris d'avoir été mis en rapport à peu près à la même époque avec un homme bien distingué dans lequel on rencontrait à un degré éminent l'alliance trop rare de la science et de la vertu. M. Teyssier né à Grenoble avait déployé dans les sciences exactes des connaissances si supérieures que quoique encore bien jeune il avait été nommé professeur à l'École polytechnique et il y jouissait de la plus haute considération. Il se servait avec zèle de l'influence qu'une position aussi élevée dans le monde savant lui avait acquise pour procurer la gloire de Dieu et la conversion des incrédules dont il était entouré. Déjà il avait obtenu des succès remarquables dans cette sorte d'apostolat et pouvait en espérer de plus grands encore lorsque le Seigneur dont les jugements sont toujours incompréhensibles ne voulut pas le laisser plus longtemps à la terre et l'appela à lui pour le mettre en possession de la récompense de sa foi et de sa charité.

M. de Noailles nous avait fait connaître en nous parlant des divers établissements qu'il avait visités l'existence d'une Congrégation à Grenoble. Cette Congrégation y avait été établie par un prêtre zélé. Elle était comme celle de Paris formée sur le modèle des anciennes Congrégations des jésuites.

% Source : PDF original, vue 071, côté gauche, page manuscrite 142.
\section*{Congrégation de Grenoble}
% Mention marginale de la page manuscrite 142.
La Congrégation de Lyon ne voulut pas demeurer en arrière dans ce commerce de charité qui s'établissait alors entre les congréganistes des différentes villes. Paris lui avait [fait] les avances ; à son tour elle en fit à Grenoble. Une association semblable à celle formée avec Paris fut la suite de cette première démarche. Grenoble nous a fourni aussi plusieurs associés ; parmi eux je vous signalerai M. de St Ferréol, chrétien modeste qui pendant un séjour assez long qu'il fit à Lyon prit part avec zèle aux travaux de nos sections.

\section*{M. de la Tour Vidaud}
% Mention marginale de la page manuscrite 142.
Ce fut encore la Congrégation de Grenoble qui nous a procuré le précieux avantage d'admettre en qualité d'associé M. de la Tour Vidaud. Plusieurs d'entre vous, Messieurs, ont connu ce fervent chrétien digne des premiers siècles de l'Église. Il vivait au milieu du monde comme un véritable anachorète ne donnant rien à aucune espèce de luxe ; toujours vêtu d'étoffes grossières, se nourrissant des aliments les plus communs, il consacrait en entier ses revenus assez considérables soit à la décoration des temples du Seigneur soit au soulagement des membres souffrants de J.-C. Il assistait parfois à nos assemblées mais bien plus souvent encore à nos fêtes. Comme nous étions heureux lorsque cette figure de prédestiné nous apparaissait dans notre petite chapelle ; nous l'apercevions modestement retiré dans quelque place obscure humblement prosterné devant la majesté suprême ; son âme aurait voulu pouvoir dérober à tous les yeux les grâces dont elle était comblée, mais malgré l'enveloppe d'humilité dans laquelle elle cherchait à se faire oublier il y avait dans
% Source : PDF original, vue 071, côté droit, page manuscrite 143.
toute sa personne un je ne sais quoi qui la décelait tout entière. Dieu permet quelquefois qu'une brillante auréole environne ainsi les saints d'ici-bas même pour que le spectacle de leurs vertus contribue à édifier leurs frères. Aussi rien ne pouvait redoubler autant la ferveur de tous les congréganistes que la présence de ce cher associé. Ne pouvaient-ils pas espérer que leurs prières seraient mieux accueillies en se trouvant en aussi bonne compagnie ? M. de Vidaud a été récemment en Savoie accusé d'avoir fait des miracles ; il est mort il y a quelques années en odeur de sainteté.

\section*{Congrégation de Chambéry}
% Mention marginale de la page manuscrite 144 ; placée au début du passage.
La Providence après avoir formé la Congrégation de Lyon semblait vouloir à l'époque dont nous nous occupons placer sous les yeux des congréganistes par le moyen des Congrégations associées des exemples de vertu de tous les genres. C'était comme une collection de modèles dont elle les enrichissait afin qu'ils pussent à loisir s'appliquer à les imiter. Ainsi après les sociétés de Paris et de Grenoble elle nous mit encore en rapport avec une autre société pieuse qui existait à Chambéry. Elle avait pour fondateur un homme dont le nom vous est devenu bien familier, M. l'abbé Rey. Il était alors secrétaire de l'évêché de Chambéry. Je ne me rappelle pas comment cette société parvint à notre connaissance mais ce que je sais parfaitement c'est que tout le bien que nous entendîmes raconter de sa société et de son directeur nous fit vivement désirer d'entrer en relation. Notre demande fut favorablement accueillie. Cette société n'avait pas pris le titre de Congrégation ; on l'appelait simplement la société des amis de Chambéry ; elle était sous le patronage du disciple de l'amour St Jean
% Source : PDF original, vue 072, côté gauche, page manuscrite 144.
l'Évangéliste. Plus éloignée de nous par ses règles que les Congrégations de Paris et de Grenoble elle s'en rapprochait davantage par son esprit ; comme nous elle cherchait à s'ensevelir dans l'obscurité d'un silence impénétrable. Parmi les membres de cette société que la Congrégation de Lyon a associés nous distinguerons soit par rapport au beau nom qu'il porte soit par rapport à ses qualités personnelles M. de Sales, héritier du nom et dernier neveu du St évêque de Genève et bien digne par ses vertus d'appartenir à ce grand saint. Lorsqu'au jour de sa réception nous l'entendîmes prononcer ces paroles aux pieds de la Ste Vierge : « Ste Marie, vierge et mère de Dieu, moi François de Sales vous choisis en ce jour pour ma souveraine, ma patronne et mon avocate... » nos cœurs furent vivement émus, nos yeux se baignèrent de larmes. Il nous semblait voir l'apôtre du Chablais en personne venu renouveller au milieu de nous sa consécration à la mère de Dieu à laquelle pendant sa vie mortelle il avait été si dévoué.

La liaison si intime qui existe depuis si longtemps entre M. Rey aujourd'hui Mgr l'évêque d'Annecy et la Congrégation commença dès lors à s'établir par correspondance. Les lettres brûlantes qu'il nous écrivait comme autant de flèches ardentes de ce cœur si chaud et si aimant venaient embraser les nôtres du feu dont il était dévoré ; mais ce ne fut que quelques années après ainsi que nous aurons à le raconter que la Providence le conduisit dans notre ville et que se conclut entre lui et nous ce traité de commerce tout spirituel qui s'est constamment exécuté depuis de part et d'autre avec la plus scrupuleuse fidélité.

% Source : PDF original, vue 072, côté droit, page manuscrite 145, avant le chapitre cinquième.
\section*{Congrégations de Bordeaux et de St Étienne}
% Mention marginale de la page manuscrite 145.
Nous nous mîmes encore en correspondance dans le même temps avec une Congrégation établie à Bordeaux à l'instar de celle de Paris. On convint du principe de l'association mais je ne crois pas que cette première démarche ait pu être suivie d'aucun résultat.

Enfin pour ne rien laisser en arrière dans ce que nous avions à vous dire de ces relations avec les Congrégations établies hors de Lyon nous ajouterons qu'on essaya dès lors d'en former une à St Étienne. Ces tentatives n'eurent presque aucun résultat. Les moments marqués par la Providence où cette société devait prendre un si grand développement et devenir pour nous un sujet d'édification n'étaient point encore arrivés ; nous aurons donc à nous en occuper plus tard.

Après avoir admiré l'action providentielle de notre bon maître pour soutenir et fortifier notre Congrégation par des moyens extérieurs cherchons à la suivre encore dans la manière dont elle favorisait son développement dans sa marche intérieure ; c'est ce qui va faire l'objet du chapitre suivant.
